\documentclass{article}
\usepackage{amsmath, amssymb, amsthm}
\usepackage{hyperref}
\usepackage{geometry}
\geometry{margin=1in}

\title{Erd\H{o}s Problem \#412}
\date{Last updated: 2025-08-31}
\author{Source: erdosproblems.com}

\begin{document}
\maketitle

\noindent\textbf{Status:} OPEN \quad \textbf{No prize}

\noindent\textbf{Tags:} number theory, iterated functions

\medskip

\noindent\textbf{Problem Statement:}

Let $\sigma_1(n)=\sigma(n)$, the sum of divisors function, and $\sigma_k(n)=\sigma(\sigma_{k-1}(n))$. Is it true that, for every $m,n\geq 2$, there exist some $i,j$ such that $\sigma_i(m)=\sigma_j(n)$?




In \cite{Er79d} Erd\H{o}s attributes this conjecture to van Wijngaarden, who told it to Erd\H{o}s in the 1950s. That is, there is (eventually) only one possible sequence that the iterated sum of divisors function can settle on. Selfridge reports numerical evidence which suggests the answer is no, but Erd\H{o}s and Graham write 'it seems unlikely that anything can be proved about this in the near future'. See also [413] and [414] .




References

[Er79d] Erd\H{o}s, P., Some unconventional problems in number theory . Acta Math. Acad. Sci. Hungar. (1979), 71-80.

\noindent\textbf{Related OEIS sequences:} A007497, A051572


\bigskip
\noindent\small{Source: \url{https://www.erdosproblems.com/412}}
\end{document}
